\documentclass[11pt]{article}
\usepackage[margin=1in]{geometry}
\usepackage{setspace}
\usepackage{enumitem}
\usepackage{amsmath}
\usepackage{amssymb}
\usepackage{fancyhdr}
\IfFileExists{lastpage.sty}{\usepackage{lastpage}}{}
\usepackage{hyperref}

\setlength{\parindent}{0pt}
\setlength{\parskip}{0.6em}
\setlist[itemize]{leftmargin=*}
\setlist[enumerate]{leftmargin=*}

\newcommand{\answerlines}[1]{\vspace{#1\baselineskip}}

\pagestyle{fancy}
\fancyhf{}
\IfFileExists{lastpage.sty}{
  \fancyfoot[C]{Financial Data Science II\ \ \ Updated March 16, 2026\ \ \ Page \thepage\ of \pageref{LastPage}}
}{
  \fancyfoot[C]{Financial Data Science II\ \ \ Updated March 16, 2026\ \ \ Page \thepage}
}

\begin{document}

\begin{center}
{\Large \textbf{Part 22: Feature Engineering}}
\end{center}

Data rarely arrive in a form that is immediately suitable for predictive
analysis. Feature engineering refers to the construction of variables from the
observed data that may be informative for a forecasting or decision problem.

In finance this step is especially important. Financial signals are weak,
relationships change over time, and even small errors in timing can invalidate
an analysis. Accordingly, the construction of features should be treated as
part of the modelling problem itself, rather than as a routine preprocessing
step.

Here we consider predictive features built from historical market data. A later
treatment could consider factor models, long-short portfolio construction, and
the relation between factor exposures and tradeable signals.

\textbf{Exercise.} Give three examples of raw data encountered in finance and
explain how each could be transformed into a useful feature for prediction.

\answerlines{6}

\newpage

\section*{Basic Terms}

A \textbf{feature} (or \textbf{predictor}) is a variable used to predict some
future quantity. A \textbf{target} (or \textbf{response}) is the quantity to be
forecast. A \textbf{forecast horizon} is the time between the observation of
the feature and the measurement of the target.

In daily equity modelling, one common target is a 5-day forward return. In
microstructure analysis, one common target is the next movement of the
mid-price. These are different problems and generally require different feature
sets.

We also distinguish between:

\begin{enumerate}
  \item \textbf{Clock time:} data aggregated to regular intervals such as one
  minute or one day.
  \item \textbf{Event time:} data indexed by market events, such as order
  submissions, cancellations, or trades.
\end{enumerate}

Finally, \textbf{lookahead bias} occurs when a feature contains information
that would not have been available at the prediction time. This is one form of
\textbf{data leakage}, and it can make a poor model look excellent in
backtesting.

\textbf{Exercise.} Suppose that you wish to predict a stock's return over the
next 5 trading days. What information is valid to use on day $t$? What
information would constitute leakage?

\answerlines{8}

\newpage

\section*{Why Feature Engineering Matters in Finance}

Feature engineering matters in any predictive problem, but several aspects of
financial data make it especially important here:

\begin{enumerate}
  \item Financial data are noisy.
  \item Relationships are often nonstationary.
  \item Timing matters. A feature that is valid at the daily level may be
  invalid at the intraday level.
  \item Useful predictors often arise from economic or market-structure
  intuition, not from generic transformations alone.
\end{enumerate}

Thus, a modest model with carefully constructed features can outperform a more
sophisticated model fitted to weak or misaligned inputs.

\section*{Example Prediction Tasks}

To fix ideas, we consider two examples in this Part.

\begin{enumerate}
  \item \textbf{Daily equity prediction:} use market and accounting data to
  predict a forward return over several days.
  \item \textbf{Limit order book prediction:} use LOBSTER message and order
  book data to predict short-horizon mid-price direction.
\end{enumerate}

The first example emphasizes feature families commonly used in equity
prediction. The second emphasizes market microstructure and the
difference between event-time data and fixed-interval bars.

\newpage

\section*{Feature Families for Daily Data}

We begin with a standard daily equity example. The objective is not to argue
that the particular features below are optimal, but to illustrate a systematic
procedure for constructing and validating daily predictors.

The feature families considered here are:

\begin{enumerate}
  \item \textbf{Returns and log returns}
  \item \textbf{Momentum over multiple horizons}
  \item \textbf{Rolling volatility}
  \item \textbf{Drawdown}
  \item \textbf{Valuation ratios such as P/E and P/B}
  \item \textbf{Cross-asset predictors such as beta, correlation, or relative momentum}
\end{enumerate}

The point of these features is not merely to create many columns. Each
predictor should encode an economically or statistically meaningful aspect of
the price process.

\section*{Feature Families for LOB Data}

We then turn to a limit order book example using LOBSTER data. These
observations are event-driven: each row corresponds to a market event together
with a synchronized snapshot of the book.

The feature families considered here are:

\begin{enumerate}
  \item \textbf{Spread and mid-price}
  \item \textbf{Level-1 imbalance}
  \item \textbf{Multi-level depth imbalance}
  \item \textbf{Order flow imbalance (OFI)}
  \item \textbf{Microprice deviation}
  \item \textbf{Message intensity and signed volume}
\end{enumerate}

These features are intended to summarize short-run pressure from demand and
supply at the book. Because the underlying data arrive in event time, one must
decide carefully how to aggregate them when building fixed-interval predictors.

\newpage

\section*{Constructing the Target}

Suppose that $P_t$ denotes a daily adjusted close price. A standard forward
return target with horizon $h$ is
\[
Y_t = \frac{P_{t+h}}{P_t} - 1.
\]

This target uses future prices by construction, but only on the right-hand side
of the modelling problem. The features used to predict $Y_t$ must be based only
on information available at time $t$.

For the LOB example, a simple directional target is
\[
Y_t =
\begin{cases}
1 & \text{if the next mid-price exceeds the current mid-price},\\
0 & \text{otherwise.}
\end{cases}
\]

\textbf{Exercise.} Explain why the following split is inappropriate for a
time-series forecasting problem:

\begin{verbatim}
from sklearn.model_selection import train_test_split
X_train, X_test, y_train, y_test = train_test_split(X, y, test_size=0.2)
\end{verbatim}

\answerlines{6}

\section*{Temporal Validation}

For the daily regression problem, one may report RMSE, MAE, and $R^2$. In
finance it is often also useful to report a directional metric, such as sign
accuracy.

For the LOB classification problem, raw accuracy can be misleading when the
classes are imbalanced. In that setting, \textbf{balanced accuracy} and
\textbf{AUC} can be more informative.

In both settings, the validation scheme should preserve temporal order. A
walk-forward or rolling-origin split is generally preferred.

\newpage

\section*{What Makes a Feature Useful?}

A useful feature need not have a large marginal correlation with the target.
More generally, a feature is attractive when it satisfies several criteria at
once:

\begin{enumerate}
  \item it carries predictive information
  \item it remains reasonably stable across time or regimes
  \item it is interpretable
  \item it can be computed in real time
  \item it survives practical constraints such as trading cost, latency, or
  limited data availability
\end{enumerate}

This is why feature engineering is better viewed as a modelling discipline than
as a list of formulas.

\section*{Exercises}

\textbf{Exercise 1.} Add rolling moving-average ratios to the daily feature set.
Re-run the daily evaluation and compare the results to those obtained from the
baseline set of return, momentum, volatility, and drawdown features.

\answerlines{6}

\textbf{Exercise 2.} In a daily equity example, explain why accounting
variables can create timing problems even when they appear to be aligned by
date.

\answerlines{6}

\textbf{Exercise 3.} In a LOBSTER example, compare level-1 imbalance and OFI.
Which is easier to interpret economically? Which appears more useful empirically?

\answerlines{6}

\textbf{Exercise 4.} Suppose that the LOBSTER data are aggregated using the
ceiling of the event time rather than the floor. Explain why this may create
leakage.

\answerlines{6}

\section*{References}

\begin{enumerate}
  \item Cont, Kukanov, and Stoikov (2013), \textit{The Price Impact of Order Book Events}.
  \item Gould, Porter, Williams, McDonald, Fenn, and Howison (2013), \textit{Limit Order Books}.
  \item LOBSTER documentation and sample files.
  \item Little and Rubin, \textit{Statistical Analysis with Missing Data}.
  \item Lopez de Prado, \textit{Advances in Financial Machine Learning}.
\end{enumerate}

\end{document}
